\Subhead{3.9.2}{Рассеяние волн на шероховатой поверхности. Понятие об обратной задаче рассеяния}

Случайная высота $h(x,y)$ меняет фазу отражённой волны. При малых наклонах и зеркальном отражении набег около $2kh\cos\theta$, где $\theta$ — угол к нормали. Для гауссовской высоты с дисперсией $\sigma_h^2$ средняя когерентная амплитуда умножается на $\exp[-2k^2\sigma_h^2\cos^2\theta]$. Когерентная интенсивность имеет квадрат этого множителя; остальная энергия может перейти в диффузное рассеяние, а не обязательно поглотиться. При больших наклонах существенны затенение и многократные отражения.

Обратная задача восстанавливает неоднородность или её статистику по измеренному полю/интенсивности. В приближении Борна разные пары направлений дают разные компоненты Фурье объекта; конечные частоты и апертуры оставляют неполное покрытие пространства $\bm q$. Измерение только интенсивности теряет фазу. Поэтому обращение обычно неединственно или неустойчиво и требует априорных ограничений, регуляризации и оценки разрешения. Реконструкция корреляционной функции случайной среды и восстановление одной конкретной реализации — разные задачи.
